% ======================================================================
%  Kapitel: Datengrundlage  --  korrigierte und gekuerzte Fassung
%
%  Kapitelfragment, keine eigene Praeambel. Vorausgesetzt werden die
%  bereits in der Arbeit definierten Makros und Umgebungen:
%    \code{...}     Monospace-Auszeichnung
%    \absatz{...}   Absatzueberschrift
%    \px            Pixel-Einheit
%    \Oe            Objektdurchmesser als Einheit
%    L{...}         linksbuendiger Spaltentyp fester Breite (array/tabularx)
%  sowie biblatex (\parencite) und die Pakete booktabs, tabularx.
%
%  Querverweise auf Kapitel 3 verwenden die Labels der korrigierten
%  Kapitelfassung: eq:deff (3.8), eq:gridmiss (3.11), eq:bendavid (3.12).
%  Der Verweis \ref{subsec:kleine-objekte} zeigt dort auf
%  \label{subsec:kleine-objekte} -- Labelnamen ggf. angleichen.
% ======================================================================

\chapter{Datengrundlage}
\label{kap:daten}

Alle Parameter der folgenden Kapitel berufen sich auf Messungen am verfügbaren
Quellmaterial. Dieses Kapitel dokumentiert sie. Sämtliche Zahlenwerte sind mit
\code{src/analyze\_data.py} an den Quelldateien gemessen und in
\code{docs/datensichtung.json} abgelegt.

\section{Annotationsdefinition}
\label{sec:annotation}

Vor jeder Messung steht die Frage, was ein Objekt ist. Aus der Sichtung der
beiden handannotierten Realbilder ergibt sich eine Definition, die die gesamte
Pipeline prägt:

\begin{quote}
Annotiert ist ausschließlich die \textbf{sichtbare Stirnfläche} eines
Rundrohrprofils -- heller verzinkter Kranz, dunkle Bohrung. Rohrschäfte in
Seitenansicht, Gitterknoten und dunkle Löcher im Maschendraht sind
\textbf{keine} Objekte, obwohl sie ähnlich aussehen. Die Box umschließt den
\textbf{sichtbaren} Teil (modale Box), nicht den geschätzten vollständigen Kreis.
\end{quote}

Beide Festlegungen sind folgenreich: Die erste bestimmt, welche Strukturen als
harte Negative benötigt werden, die zweite, wie die Boxen im Generator berechnet
werden. Die Unterscheidung modal / amodal ist in der Literatur etabliert
\parencite{li2016amodal,zhu2017amodal}; amodale Boxen umschließen auch verdeckte
Teile. Erzeugte der Generator amodale Boxen, während die Realbilder modal
annotiert sind, entstünde ein systematischer Größenversatz -- ein Verstoß gegen
Anforderung~A5.

\section{Quellmaterial}
\label{sec:quellmaterial}

\absatz{Objektausschnitte.} Verfügbar sind neun aus Realaufnahmen
ausgeschnittene Rohrprofile der Größe $28\times30$ bis $43\times38$ Pixel. Sechs
zeigen eine dunkle Bohrung (offenes Rohr), drei eine helle (Sonnenlicht fällt in
die Röhre oder das Ende ist gestaucht); beide Formen kommen in den Realbildern
vor, weshalb gleichverteilt aus allen neun gezogen wird. Die aus dem
Ellipsenmodell gewonnenen Achsverhältnisse der Ausschnitte streuen zwischen 0,83
und 1,14 (Abschnitt~\ref{sec:matting}). Kritisch ist Befund~1 in
Tabelle~\ref{tab:befunde}: Alle neun sind zwar RGBA-Dateien, ihr Alphakanal ist
aber durchgängig 255 -- es handelt sich um rechteckige Ausschnitte, bei denen
13--35\,\% der Fläche nicht zum Objekt gehören.

\absatz{Hintergründe.} Von den sieben Hintergrundaufnahmen sind vier
\emph{Clean Plates}: Wände aus Rohrschäften derselben Objektart \emph{ohne}
Stirnflächen, mit den gemessenen Rohrdurchmessern 16/18/21/20\px{}. Sie liefern
gleichzeitig photometrisch passenden Kontext, das Rohmaterial der
Okkluder-Bibliothek und echte Negativbeispiele und werden deshalb bevorzugt
gezogen: In 80\,\% der Bilder stammt der Hintergrund aus dieser Gruppe, was je
Clean Plate 20\,\% und je Fremdszene 6,7\,\% ergibt -- ein dreifaches Gewicht je
Bild. In der Terminologie von Gleichung~\eqref{eq:bendavid} verkleinern sie
denjenigen Anteil des Divergenzterms, der auf den Hintergrund zurückgeht: Dieser
ist nicht \emph{ähnlich} der Zieldomäne, er stammt aus ihr. Die Komposition
selbst bleibt davon unberührt und ist weiterhin eine Divergenzquelle
(Signaturen~S1--S6).

Die drei Fremdszenen (Kieswerk, Betonplatten, Betondeck mit Maschendrahtzaun)
sind primär als hartes Negativmaterial wertvoll -- der Maschendraht ist genau
die Struktur, die nach der Annotationsdefinition kein Objekt ist und dennoch als
Falschpositiv droht. Schwäche des Materials: Alle vier Clean Plates zeigen
dieselbe Stapelart bei bedecktem bis diffusem Licht -- kein harter
Sonnenschein, keine Nässe, kein Schnee.

\section{Bewertungsdaten und ihre Statistik}
\label{sec:valtest}

Zur Bewertung stehen zwei vollständig handannotierte Realaufnahmen zur
Verfügung: \code{val} ($2459\times1174$, 98 Objekte) für die Modellauswahl und
\code{test} ($1600\times900$, 89 Objekte) ausschließlich für die Endbewertung.
Beide Annotationen wurden nach der ersten Auswertungsrunde noch einmal
durchgesehen und ergänzt; alle berichteten Zahlen beziehen sich auf den so
korrigierten Stand mit zusammen 187 Objekten.

\begin{table}[htbp]
\centering
\caption[Gemessene Statistik der realen Labels]{Gemessene Statistik der realen
Labels ($n = 187$). Als Objektdurchmesser dient $\sqrt{\text{Fläche}}$; alle
Abstände sind in Objektdurchmessern (\Oe) angegeben und damit skaleninvariant
direkt als Generatorparameter verwendbar. NN = nächster Nachbar.}
\label{tab:statistik}
\footnotesize
\begin{tabular}{@{}lll@{}}
\toprule
\textbf{Größe} & \textbf{\code{val}} & \textbf{\code{test}} \\
\midrule
Anzahl Objekte & 98 & 89 \\
Durchmesser (px) & 12,2--35,8, \textbf{Median 26,8} & 6,4--31,9, \textbf{Median 17,0} \\
\quad p05 / p90 & 19,1 / 33,0 & 12,2 / 27,4 \\
Seitenverhältnis $w/h$ & p10 0,97, Median 1,09, p90 1,40 & p10 0,82, Median 1,08, p90 1,22 \\
Objektfläche gesamt & 2,60\,\% des Bildes & 2,65\,\% \\
am Bildrand angeschnitten & 2,0\,\% & 6,7\,\% \\
\midrule
NN-Abstand (in \Oe) & p10 0,86, Median 1,11, p90 1,78 & p10 0,78, Median 1,26, p90 2,72 \\
Nachbarn innerhalb $5\,\Oe$ & Median 10, max 15 & Median 10, max 20 \\
NN näher als $1{,}2\,\Oe$ & 60,2\,\% & 44,9\,\% \\
Größenverhältnis zum NN & Median 1,05, p90 1,18 & Median 1,07, p90 1,38 \\
Nachbarrichtung & quasi isotrop, Häufung $90$--$150^\circ$ & quasi isotrop \\
\bottomrule
\end{tabular}
\end{table}

Mit Durchmessern von 6--36\px{} liegen die Objekte überwiegend unterhalb der
COCO-Grenze für Kleinobjekte von $32^2$\px{} Fläche \parencite{lin2014coco};
lediglich das obere Ende von \code{val} überschreitet sie. Nach
Gleichung~\eqref{eq:deff} schrumpfen die Mediane beider Splits bei
\code{imgsz}=640 auf etwa 7\px, die Hauptmasse zwischen p05 und p90 auf
5--11\px. Damit liegt praktisch die gesamte Objektmasse in der Größenordnung des
P3-Strides von 8\px{} oder darunter, und der in
Gleichung~\eqref{eq:gridmiss} beschriebene Fall wird für den Großteil der
Instanzen geometrisch möglich. Die Inferenzauflösung wird deshalb in
Abschnitt~\ref{sec:imgsz} systematisch untersucht.

Zwei Zahlen der Tabelle sind für die weiteren Kapitel besonders wichtig.
Erstens sind vier der 89 Objekte in \code{test} schon im \emph{nativen} Maßstab
kleiner als 12\px, zwei davon kleiner als die kritische Größe $S\sqrt{2}$ aus
Gleichung~\eqref{eq:gridmiss}, die für den P3-Stride $S=8$ der Werkzeugkette
\parencite{ultralyticsdocs} bei 11,3\px{} liegt. Für
diese Objekte existieren nach der idealisierten Betrachtung in
Abschnitt~\ref{subsec:kleine-objekte} Gitterlagen, in denen kein Zentrum der
P3-Ebene innerhalb des Objekts liegt. Das ist keine Detektionsgrenze, aber ein
konkreter Anlass für die Modellwahl in Abschnitt~\ref{sec:p2kopf}. Zweitens sind
die beiden Bilder \emph{nicht} austauschbar: \code{test} enthält kleinere
Objekte (Median 17,0 gegen 26,8\px), lockerere Packung (NN-Median 1,26 gegen
1,11\,\Oe) und mehr angeschnittene Boxen (6,7 gegen 2,0\,\%). Der schwierigere
Split ist damit gerade derjenige, den Training und Modellauswahl nie berühren --
was bei der Lektüre der Ergebnisse zu berücksichtigen ist.

\section{Befunde und die daraus abgeleiteten Parameter}
\label{sec:messungmapping}

Tabelle~\ref{tab:befunde} fasst zusammen, welcher Befund welchen
Pipelineparameter festlegt. Sie ist die Belegbrücke zu
Kapitel~\ref{kap:pipeline} und macht zugleich sichtbar, welche Parameter
\emph{nicht} durch eine Messung gedeckt sind. Die Spalte \emph{Art} bezieht sich
dabei auf die Begründung des \emph{abgeleiteten Parameters}, nicht auf die des
Befunds; eine Doppelangabe wie \emph{M/A} bedeutet, dass der Befund gemessen,
der konkrete Zahlenwert des Parameters aber gewählt ist.

\begin{table}[htbp]
\centering
\caption[Befunde und abgeleitete Parameter]{Befunde der Datensichtung und die
daraus abgeleiteten Parameter. Belegart: \emph{M} = Messung an den Realdaten,
\emph{P} = physikalische Randbedingung, \emph{L} = Literatur, \emph{A} =
Annahme.}
\label{tab:befunde}
\footnotesize
\begin{tabularx}{\textwidth}{@{}L{0.30\textwidth}XL{0.05\textwidth}@{}}
\toprule
\textbf{Befund / Messwert} & \textbf{abgeleiteter Parameter} & \textbf{Art} \\
\midrule
\textbf{1} Alphakanal durchgängig 255; 13--35\,\% Fremdkontext &
eigene Matting-Stufe mit weicher Alphamaske (Signatur~S1) & M \\
\textbf{2} sechs dunkle, drei helle Bohrungen; Achsverhältnis der Ausschnitte
0,83--1,14 &
Gleichverteilung über alle neun Ausschnitte; Ovalisierung überhaupt vorsehen & M/A \\
\textbf{3} vier Clean Plates der Zieldomäne; Maschendraht und Schotter in den
Fremdszenen &
80\,\% Ziehgewicht auf Clean Plates; 10\,\% Negativbilder (Gewichte gewählt) & M/A \\
\textbf{4} Rohrdurchmesser der Clean Plates 16--21\px{} &
Kopplung von Zoom und Objektgröße & M/P \\
\textbf{5} Objektdurchmesser 6,4--35,8\px, p05 der beiden Splits 19,1 bzw.\
12,2\px{} &
Generatorbereich 12--44\px: untere Grenze $=$ p05 von \code{test}, obere
$=$ Maximum $+\,25\,\%$ Reserve (Faktor gewählt); die vier Objekte unter
12\px{} in \code{test} deckt erst die Trainingsskalierung ab & M/A \\
\textbf{6} Seitenverhältnis p10--p90 $=$ 0,82--1,40 &
Ovalisierung bei 57\,\% der Instanzen, oblat:prolat $=$ 45:12; Quote als
Anpassung an das p10--p90-Band gewählt & M/A \\
\textbf{7} Größenverhältnis zum Nachbarn: Median 1,05, p90 1,18 (\code{val});
Median 1,07, p90 1,38 (\code{test}) &
eine Basisgröße pro Cluster, je Instanz nur $\pm 8\,\%$; abgeleitet aus dem p90
von \code{val}, das p90 von \code{test} wird damit nicht abgedeckt & M/A \\
\textbf{8} hinter jeder Stirnfläche führt ein Rohrschaft weg, ist aber nicht
immer sichtbar &
\emph{sichtbarer} Schaft bei 88\,\% der Instanzen (Tabelle~\ref{tab:parameter},
T5), Anteil gewählt; Schaftmaterial aus der photographischen
Okkluder-Bibliothek & P/A \\
\textbf{9} NN-Abstand Median 1,11--1,26\,\Oe; 6,7\,\% angeschnittene Labels in
\code{test} &
Gitterabstände des Layoutmodells; Clusterzentrum darf außerhalb der Leinwand
liegen, Randregel 15\,\% (Wert gewählt) & M/A \\
\midrule
\textbf{10} keine Messung verfügbar &
Anteil verformter Profile (18\,\%), Schrägansicht (30\,\%), Schlagschatten
(49\,\%) & A \\
\textbf{11} je \emph{ein} Bild pro Split &
Bootstrap-Konfidenzintervalle statt Punktschätzungen
(Abschnitt~\ref{sec:bootstrap}) & M \\
\bottomrule
\end{tabularx}
\end{table}

Befund~11 ist die größte methodische Schwäche der Arbeit: Je ein Bild pro Split
ist für eine belastbare Metrikangabe zu wenig, und beide Bilder zeigen dieselbe
Objektart bei ähnlichem Wetter. Alle berichteten Zahlen sind entsprechend als
Punktschätzungen mit erheblicher Unsicherheit zu lesen.
